% 贪心近似界的来源：每轮差距至少缩小 1/k（传感器覆盖实例，k=2，最优集合 O={b,c}，OPT=16）。
% 第1轮：S0=空集，F=0，差距16；O 中元素相对 S0 的增益 b=8、c=8，合计16>=差距，最好者>=16/2=8；贪心选 a，增益10。
% 第2轮：S1={a}，F=10，差距6；O 中元素相对 S1 的增益 b=2、c=4，合计6>=差距，最好者>=6/2=3；贪心选 c，增益4，差距降为2。
\begin{tikzpicture}[x=4mm,y=1mm,line width=.6pt,font=\figurefont\fontsize{8}{10}\selectfont]
  \tikzset{lab/.style={FigureInk,anchor=east,font=\figurefont\fontsize{7.5}{9}\selectfont},
           seg/.style={FigureInk,font=\figurefont\fontsize{7.5}{9}\selectfont},
           kmark/.style={FigureRed,dashed,line width=.8pt},
           klab/.style={FigureRed,anchor=east,font=\figurefont\fontsize{7.5}{9}\selectfont}}
  % 坐标参考线
  \draw[FigureStroke,line width=.75pt] (0,-49)--(0,4);
  \node[FigureInk,anchor=south] at (0,4.2) {$0$};
  \draw[FigureStroke,line width=.75pt] (16,-49)--(16,4);
  \node[FigureInk,anchor=south] at (16,4.2) {$\mathrm{OPT}=16$};
  % ---------- 第1轮 ----------
  \node[FigureInk,anchor=west] at (-6.6,7.5) {第1轮：$S_0=\varnothing$};
  \node[lab] at (-0.25,0) {当前值与差距};
  \draw[FigureMuted,dashed] (0,-2) rectangle (16,2);
  \node[seg] at (12,0) {差距$16$};
  \node[lab] at (-0.25,-7) {$O$中元素的增益};
  \filldraw[fill=FigureGreenFill,draw=FigureGreen] (0,-9) rectangle (8,-5);
  \node[seg] at (4,-7) {$b{:}\,8$};
  \filldraw[fill=FigureGreenFill,draw=FigureGreen] (8,-9) rectangle (16,-5);
  \node[seg] at (12,-7) {$c{:}\,8$};
  \node[lab] at (-0.25,-14) {贪心选中};
  \filldraw[fill=FigureYellowFill,draw=FigureYellow] (0,-16) rectangle (10,-12);
  \node[seg] at (5,-14) {$a{:}\,10$};
  \draw[kmark] (8,-17.5)--(8,3.5);
  \node[klab] at (7.9,4.6) {差距的$1/k$：$8$};
  % ---------- 第2轮 ----------
  \begin{scope}[yshift=-31mm]
  \node[FigureInk,anchor=west] at (-6.6,7.5) {第2轮：$S_1=\{a\}$};
  \node[lab] at (-0.25,0) {当前值与差距};
  \filldraw[fill=FigureBlueFill,draw=FigureBlue] (0,-2) rectangle (10,2);
  \node[seg] at (5,0) {$F(S_1)=10$};
  \draw[FigureMuted,dashed] (10,-2) rectangle (16,2);
  \node[seg] at (14.5,0) {差距$6$};
  \node[lab] at (-0.25,-7) {$O$中元素的增益};
  \filldraw[fill=FigureGreenFill,draw=FigureGreen] (10,-9) rectangle (12,-5);
  \node[seg] at (11,-7) {$b{:}\,2$};
  \filldraw[fill=FigureGreenFill,draw=FigureGreen] (12,-9) rectangle (16,-5);
  \node[seg] at (14.6,-7) {$c{:}\,4$};
  \node[lab] at (-0.25,-14) {贪心选中};
  \filldraw[fill=FigureYellowFill,draw=FigureYellow] (10,-16) rectangle (14,-12);
  \node[seg] at (11.5,-14) {$c{:}\,4$};
  \draw[kmark] (13,-17.5)--(13,3.5);
  \node[klab] at (12.9,4.6) {差距的$1/k$：$3$};
  \end{scope}
\end{tikzpicture}
